\section{FORMAL RESTATEMENT}
\textbf{Erd\H{o}s \#67 (the Erd\H{o}s discrepancy problem); partial Lean 4 formalization.}
For $f:\mathbb{N}\to\{-1,+1\}$ and every real $C>0$, there exist $d,m\ge1$ with
\[
 \Bigl|\sum_{k=1}^{m} f(kd)\Bigr| > C .
\]
This is \texttt{Erdos67.erdos\_67} in formal-conjectures, with $f$ valued in the finset $\{-1,1\}\subset\mathbb{R}$. The complex variant \texttt{Erdos67.erdos\_67.variants.complex} takes $f:\mathbb{N}\to S^1\subset\mathbb{C}$ and the complex norm.

\section{QUICK LITERATURE/CONTEXT CHECK}
Both statements were proved by T.~Tao, \emph{The Erd\H{o}s discrepancy problem}, Discrete Analysis 2016:1, \url{https://arxiv.org/abs/1509.05363}. See \url{https://www.erdosproblems.com/67} and the formal statement \url{https://github.com/google-deepmind/formal-conjectures/blob/main/FormalConjectures/ErdosProblems/67.lean}. (The reference list of that file links to \texttt{erdosproblems.com/66}; the intended page is \texttt{/67}.) Tao's proof relies on logarithmically averaged correlation estimates for multiplicative functions. It is not formalized here, and no novelty is claimed.

The unconditional part below is the classical fact that every $\pm1$ sequence has discrepancy at least $2$: a sequence with all homogeneous partial sums in $\{-1,0,1\}$ has length at most $11$.

\section{ATTACK PLAN}
\begin{enumerate}
\item Prove the real statement for every $C<2$ by showing that $|\sum_{k\le m}f(kd)|\le C<2$ for all $d,m$ is contradictory. Only the values $f(1),\dots,f(12)$ are needed.
\item Prove that the complex variant implies the real statement, since $\pm1\in S^1$ and the complex norm of a real number is its absolute value.
\end{enumerate}

\section{WORK}
Write $x_n=f(n)$. If two signs satisfy $|a+b|<2$ then $b=-a$. Assume every homogeneous partial sum has absolute value at most $C<2$.
\begin{itemize}
\item $d=1$, $m=2,4,6,10$: successively $x_2=-x_1$, $x_4=-x_3$, $x_6=-x_5$, $x_8=-x_7$, $x_{10}=-x_9$ (each even partial sum reduces to its last two terms).
\item $d=2$, $m=2,4,6$: $x_4=-x_2$, $x_8=-x_6$, $x_{12}=-x_{10}$.
\item $d=3$, $m=2,4$: $x_6=-x_3$, $x_{12}=-x_9$.
\item $d=5$, $m=2$: $x_{10}=-x_5$.
\end{itemize}
Then $x_{12}=-x_{10}=x_5=-x_6=x_3$, while $x_{12}=-x_9=x_{10}=-x_5=x_6=-x_3$. So $x_3=-x_3$, impossible for $x_3=\pm1$.

Full Lean 4 source (\texttt{Erdos67Attempt.lean}):
\begin{verbatim}
import Mathlib.Analysis.Complex.Norm
import Mathlib.Data.Complex.BigOperators
import Mathlib.Order.Interval.Finset.Nat
import Mathlib.Algebra.Order.BigOperators.Group.Finset

/-!
# Erdős problem 67 (the Erdős discrepancy problem): partial formalization

STATUS: PARTIAL. Tao's theorem [Ta16] is NOT formalized here. We prove, with no `sorry`
and no custom axioms:

* `erdos_67_of_lt_two`: the formal-conjectures statement `Erdos67.erdos_67` holds
  unconditionally for every `0 < C < 2`, i.e. every `±1` sequence has discrepancy at least `2`.
  Only the values `f 1, …, f 12` and the differences `d = 1, 2, 3, 5` are used.
* `erdos_67_of_complex`: the complex variant `Erdos67.erdos_67.variants.complex` implies
  the real statement `Erdos67.erdos_67`, since `{-1, 1} ⊆ S¹`.
-/

namespace Erdos67Attempt

/-- The two admissible values. -/
theorem val_cases {x : ℝ} (hx : x ∈ ({-1, 1} : Finset ℝ)) : x = -1 ∨ x = 1 := by
  simpa using hx

/-- Two signs whose sum has absolute value below `2` are opposite. -/
theorem neg_of_abs_add_lt_two {a b : ℝ} (ha : a = -1 ∨ a = 1) (hb : b = -1 ∨ b = 1)
    (h : |a + b| < 2) : b = -a := by
  rcases ha with rfl | rfl <;> rcases hb with rfl | rfl
  · rw [show (-1 : ℝ) + -1 = -2 by norm_num, abs_neg, abs_two] at h
    exact absurd h (lt_irrefl _)
  · norm_num
  · norm_num
  · rw [show (1 : ℝ) + 1 = 2 by norm_num, abs_two] at h
    exact absurd h (lt_irrefl _)

/-- **Discrepancy at least 2.** For every `C < 2` the formal-conjectures statement holds. -/
theorem erdos_67_of_lt_two (f : ℕ → ({-1, 1} : Finset ℝ)) (C : ℝ) (hC : C < 2) :
    ∃ᵉ (d ≥ 1) (m ≥ 1), C < |∑ k ∈ Finset.Icc 1 m, (f (k * d)).1| := by
  by_contra hno
  push Not at hno
  set x : ℕ → ℝ := fun n => (f n).1 with hx
  have hv : ∀ n, x n = -1 ∨ x n = 1 := fun n => val_cases (f n).2
  have key : ∀ d m, 1 ≤ d → 1 ≤ m → |∑ k ∈ Finset.Icc 1 m, x (k * d)| < 2 :=
    fun d m hd hm => lt_of_le_of_lt (hno d hd m hm) hC
  -- the partial sums we need, written out
  have s12 : |x 1 + x 2| < 2 := by
    simpa [Finset.sum_Icc_succ_top] using key 1 2 le_rfl (by norm_num)
  have s14 : |x 1 + x 2 + x 3 + x 4| < 2 := by
    simpa [Finset.sum_Icc_succ_top] using key 1 4 le_rfl (by norm_num)
  have s16 : |x 1 + x 2 + x 3 + x 4 + x 5 + x 6| < 2 := by
    simpa [Finset.sum_Icc_succ_top] using key 1 6 le_rfl (by norm_num)
  have s110 : |x 1 + x 2 + x 3 + x 4 + x 5 + x 6 + x 7 + x 8 + x 9 + x 10| < 2 := by
    simpa [Finset.sum_Icc_succ_top] using key 1 10 le_rfl (by norm_num)
  have s18 : |x 1 + x 2 + x 3 + x 4 + x 5 + x 6 + x 7 + x 8| < 2 := by
    simpa [Finset.sum_Icc_succ_top] using key 1 8 le_rfl (by norm_num)
  have s22 : |x 2 + x 4| < 2 := by
    simpa [Finset.sum_Icc_succ_top] using key 2 2 (by norm_num) (by norm_num)
  have s26 : |x 2 + x 4 + x 6 + x 8 + x 10 + x 12| < 2 := by
    simpa [Finset.sum_Icc_succ_top] using key 2 6 (by norm_num) (by norm_num)
  have s24 : |x 2 + x 4 + x 6 + x 8| < 2 := by
    simpa [Finset.sum_Icc_succ_top] using key 2 4 (by norm_num) (by norm_num)
  have s32 : |x 3 + x 6| < 2 := by
    simpa [Finset.sum_Icc_succ_top] using key 3 2 (by norm_num) (by norm_num)
  have s34 : |x 3 + x 6 + x 9 + x 12| < 2 := by
    simpa [Finset.sum_Icc_succ_top] using key 3 4 (by norm_num) (by norm_num)
  have s52 : |x 5 + x 10| < 2 := by
    simpa [Finset.sum_Icc_succ_top] using key 5 2 (by norm_num) (by norm_num)
  -- propagate the forced signs
  have e2 : x 2 = -x 1 := neg_of_abs_add_lt_two (hv 1) (hv 2) s12
  have e4 : x 4 = -x 3 := neg_of_abs_add_lt_two (hv 3) (hv 4)
    (by have h : x 3 + x 4 = x 1 + x 2 + x 3 + x 4 := by rw [e2]; ring
        rw [h]; exact s14)
  have e6 : x 6 = -x 5 := neg_of_abs_add_lt_two (hv 5) (hv 6)
    (by have h : x 5 + x 6 = x 1 + x 2 + x 3 + x 4 + x 5 + x 6 := by rw [e2, e4]; ring
        rw [h]; exact s16)
  have e8 : x 8 = -x 7 := neg_of_abs_add_lt_two (hv 7) (hv 8)
    (by have h : x 7 + x 8 = x 1 + x 2 + x 3 + x 4 + x 5 + x 6 + x 7 + x 8 := by
          rw [e2, e4, e6]; ring
        rw [h]; exact s18)
  have e10 : x 10 = -x 9 := neg_of_abs_add_lt_two (hv 9) (hv 10)
    (by have h : x 9 + x 10 =
            x 1 + x 2 + x 3 + x 4 + x 5 + x 6 + x 7 + x 8 + x 9 + x 10 := by
          rw [e2, e4, e6, e8]; ring
        rw [h]; exact s110)
  have f4 : x 4 = -x 2 := neg_of_abs_add_lt_two (hv 2) (hv 4) s22
  have f8 : x 8 = -x 6 := neg_of_abs_add_lt_two (hv 6) (hv 8)
    (by have h : x 6 + x 8 = x 2 + x 4 + x 6 + x 8 := by rw [f4]; ring
        rw [h]; exact s24)
  have f12 : x 12 = -x 10 := neg_of_abs_add_lt_two (hv 10) (hv 12)
    (by have h : x 10 + x 12 = x 2 + x 4 + x 6 + x 8 + x 10 + x 12 := by
          rw [f4, f8]; ring
        rw [h]; exact s26)
  have g6 : x 6 = -x 3 := neg_of_abs_add_lt_two (hv 3) (hv 6) s32
  have g12 : x 12 = -x 9 := neg_of_abs_add_lt_two (hv 9) (hv 12)
    (by have h : x 9 + x 12 = x 3 + x 6 + x 9 + x 12 := by rw [g6]; ring
        rw [h]; exact s34)
  have h10 : x 10 = -x 5 := neg_of_abs_add_lt_two (hv 5) (hv 10) s52
  -- x 12 = -x 10 = x 5 = -x 6 = x 3 and x 12 = -x 9 = x 10 = -x 5 = x 6 = -x 3
  have hA : x 12 = x 3 := by linarith
  have hB : x 12 = -x 3 := by linarith
  rcases hv 3 with h3 | h3 <;> linarith

/-- The complex variant implies the real statement. -/
theorem erdos_67_of_complex
    (hcx : ∀ (g : ℕ → Metric.sphere (0 : ℂ) 1) (C : ℝ), 0 < C →
      ∃ᵉ (d ≥ 1) (m ≥ 1), C < ‖∑ k ∈ Finset.Icc 1 m, (g (k * d)).1‖)
    (f : ℕ → ({-1, 1} : Finset ℝ)) (C : ℝ) (hC : 0 < C) :
    ∃ᵉ (d ≥ 1) (m ≥ 1), C < |∑ k ∈ Finset.Icc 1 m, (f (k * d)).1| := by
  let g : ℕ → Metric.sphere (0 : ℂ) 1 := fun n =>
    ⟨((f n).1 : ℂ), by
      rw [mem_sphere_zero_iff_norm, Complex.norm_real, Real.norm_eq_abs]
      rcases val_cases (f n).2 with h | h <;> rw [h] <;> norm_num⟩
  obtain ⟨d, hd, m, hm, hlt⟩ := hcx g C hC
  refine ⟨d, hd, m, hm, ?_⟩
  have : (∑ k ∈ Finset.Icc 1 m, (g (k * d)).1) =
      ((∑ k ∈ Finset.Icc 1 m, (f (k * d)).1 : ℝ) : ℂ) := by
    rw [Complex.ofReal_sum]
  rw [this, Complex.norm_real, Real.norm_eq_abs] at hlt
  exact hlt

#print axioms erdos_67_of_lt_two
#print axioms erdos_67_of_complex

end Erdos67Attempt
\end{verbatim}

\section{VERIFICATION}
Checked with \texttt{lake env lean Erdos67Attempt.lean} using Lean v4.33.1 and Mathlib revision \texttt{0df444a360eaa60ab8c11dca51a86af692955474} (the revision pinned by formal-conjectures). No errors, no warnings, no \texttt{sorry}. \texttt{\#print axioms} reports only \texttt{propext}, \texttt{Classical.choice}, \texttt{Quot.sound} for both theorems. The conclusions use exactly the expressions of the formal-conjectures statements. \texttt{erdos\_67\_of\_lt\_two} does not even need $C>0$.

\section{FINAL STATUS}
\textbf{PARTIAL.} Proved in Lean: \texttt{erdos\_67} for all $C<2$, unconditionally; and \texttt{erdos\_67.variants.complex} $\Rightarrow$ \texttt{erdos\_67}. Not formalized: Tao's theorem for $C\ge2$. The first remaining gap is $C=2$, where the known elementary route is a large computer search. The general case needs Tao's analytic argument. Subject to expert review.
